\section{Introduction}
Pretraining can improve its objective while making a representation less useful for a later task. The difficult question is structural: which updates change task readability, and what determines the size and duration of that change? Output variance, spectral concentration, and total gradient norm do not answer this question, because unregularized population linear risk depends on a subspace rather than on its coordinate scaling.

We study a convolutional encoder with squaring, global average pooling, and a linear head. Its local second-moment representation is exact for arbitrary image distributions: each channel reads $h_c^\top M(x)h_c$. Whitening the centered moment dictionary puts coordinate tasks, band-energy tasks, and arbitrary multi-output labels in one regression framework. The risk formula is standard least squares; its role is to identify the task matrix against which subspace motion must be evaluated. Independent Gaussian frequencies are a tractable special case, rather than a prerequisite for representing natural images.

The first structural distinction concerns the source of exterior motion. A linear factorization cannot move its encoder row space under an isotropic population objective. Away from isotropy, we bound its normal velocity by an explicit residual gradient. A square encoder has an additional source: anisotropy of its channel Jacobian Gram operators. For Gaussian patches these operators are Jordan products of rank-one filter matrices. With two independent filters, encoder motion requires both nonzero filter correlation and objective/head cross coupling. This criterion also supplies nonlinear cases with exact conservation.

We use that structure to solve an actual parameter flow with $K=1<C=2<d=3$. A moving quadratic direction passes through task-aligned directions and continues beyond them. An open cone of tasks shares two comparison states at which both encoder and projection risks have improved and then exceeded initialization, with explicit uniform gaps. The family covers a continuum of scales, orientations, learning-rate ratios, and a structured five-percent covariance range. Its symmetry is a checkable sufficient initialization condition, not a high-probability statement. Within this symmetry, different rates only reparameterize the ordered subspace path, so rebound alone does not imply rate competition.

For general states we bound finite risk change by initial task cross-block mass and normal subspace movement. This separates small motion from weak task coupling as two amplitude bottlenecks. Two nonzero block projector tangents admit a competing task exactly when they are not positively proportional. Finite-budget and error bounds then specify how state drift, gradient uncertainty, and step size consume a source competition margin. These general bounds may be conservative; they are not useful merely because their constants exist. On pre-existing matched-function population trajectories, source-only interval boxes certify finite risk changes of opposite signs at a common energy section, and the saved future values fall within their amplitude envelopes. This is a retrospective nonvacuity check with prediction constants separated from future records.

The original computer-assisted fixed-gap projection construction and controlled texture experiments are retained. They distinguish selection inside an encoder space from changes to that space, and expose systematic finite-data linear drift and late-stage prediction failure. Natural-image evidence is compressed to a boundary table, with full diagnostics in the appendix. CIFAR head-selection gaps largely predate training, and suppressing a dominant DTD statistic does not restore class learning. The new task coupling, normal budget, and error constants have not been measured there; we do not claim to explain those failures. The resulting contribution is a structural analysis with explicit constructive and numerical domains, rather than a general claim that pretraining or a particular rate ratio loses semantic information.

\section{Related Work}
\paragraph{Initialization and optimization geometry.}
Kunin et al.~\cite{kunin2024rich} analyze how layer imbalance and learning rates change feature learning, including comparisons with the same initial product. Our scale--learning-rate equivalence belongs to this broader setting. Its role here is to connect the actual parameter metric to an explicit independent-task risk threshold, rather than to establish scale sensitivity for the first time.

\paragraph{Projection heads and feature selection.}
Gui et al.~\cite{gui2023projection} characterize expansion and shrinkage induced by projection heads with fixed encoders. Xue et al.~\cite{xue2024benefits} study layerwise feature weighting and report settings with post-projection early-stopping gains. Torralba and Weiss~\cite{torralba2026theory} analyze the same broad convolution--nonlinearity--GAP architecture using GUPA; their Appendix M discusses nonunique decompositions, initialization effects, and recognition performance falling while the loss improves. Our question concerns finite-time population risk along actual joint parameter training, with fixed risk gaps and explicit encoder/projection separation.

\paragraph{Nonlinearity and collapse.}
Jing et al.~\cite{jing2022collapse} analyze two-layer linear contrastive dynamics and dimensional collapse; Tian~\cite{tian2023nonlinearity} studies the role of nonlinearity in contrastive dynamics. Chaudhry~\cite{chaudhry2026geometry} develops a geometric account of projection heads. The Jacobian Gram matrix governing our effective updates is distinct from a pullback loss Hessian. Moreover, our scalar variance remains positive and our texture experiments retain their numerical ranks: changing a task direction is not, by itself, dimensional collapse. Our contribution is scoped to the stated construction and measured models, without an exhaustive priority claim over all equivalent factorizations.
